\documentclass[11pt]{amsart}
\usepackage[T1]{fontenc}
\usepackage{lmodern,amsmath,amssymb,amsthm,microtype}
\usepackage[margin=1.15in]{geometry}
\usepackage[colorlinks=true,linkcolor=blue,citecolor=blue,urlcolor=blue]{hyperref}
% Repository locator for the public (non-anonymous) build of
% non_mf_groups_exist.tex.  Kept in a separate file so that the anonymous
% build has an anonymous source archive and not merely an anonymous PDF:
% exclude this file from a double-blind submission and the manuscript still
% compiles, printing "supplied with this submission as supplementary
% material" in place of the link.
%
% Loaded only when \anonymousfalse is set in the preamble.
\newcommand{\coderepository}{https://github.com/SauersML/group-approximation}
\newcommand{\coderepositoryname}{github.com/SauersML/group-approximation}

\hypersetup{pdftitle={An infinite simple Kazhdan sofic group},pdfauthor={\paperauthorname}}
\newtheorem{theorem}{Theorem}
\newtheorem{corollary}[theorem]{Corollary}
\newcommand{\F}{\mathbb F}
\newcommand{\Z}{\mathbb Z}
\newcommand{\LC}{\operatorname{LC}}
\newcommand{\EL}{\operatorname{EL}}
\newcommand{\GL}{\operatorname{GL}}
\newcommand{\PSL}{\operatorname{PSL}}
\newcommand{\diag}{\operatorname{diag}}
\newcommand{\doi}[1]{\href{https://doi.org/#1}{\nolinkurl{doi:#1}}}
\raggedbottom
\emergencystretch=1.5em

\begin{document}
\title{An infinite simple Kazhdan sofic group}
\author{\paperauthorname}
\date{}
\subjclass[2020]{Primary 20E32; Secondary 20E26, 20F10, 22D10, 37B10, 16S35}

\begin{abstract}
For every infinite minimal subshift $X$, the group
$\EL_3(\LC(X,\F_2)\rtimes\Z)$ is infinite, simple, finitely generated,
and has property~\textup{(T)}. It is locally embeddable into finite
groups, so it is sofic and hyperlinear. This answers the question of
Brown and Ozawa whether an infinite simple Kazhdan group can be
hyperlinear, and Pestov's sofic version of it. Every finitely generated
residually finite group is a subgroup of an infinite simple Kazhdan LEF
group, and every Turing degree occurs as the word-problem degree of one
of these groups.
\end{abstract}
\maketitle

Can an infinite simple Kazhdan group be hyperlinear? Brown asked this
in its von Neumann algebra form in 2001~\cite[\S11, Question~7]{Brown},
and Ozawa in the hyperlinear form in 2003~\cite[p.~527]{Ozawa}.
Pestov's Open question~9.1 adds the sofic version~\cite{Pestov}.
The following construction answers all three forms positively.

\begin{theorem}\label{thm:main}
Let $X\subseteq A^{\Z}$ be an infinite minimal subshift over a finite
alphabet, with shift $T$, and let $\LC(X,\F_2)$ be the ring of locally
constant functions $X\to\F_2$. Then
\[
  G_X=\EL_3\bigl(\LC(X,\F_2)\rtimes_T\Z\bigr)
\]
is an infinite, finitely generated, simple group with Kazhdan's
property~\textup{(T)}. It is locally embeddable into finite groups
\textup{(LEF)}, so it is sofic and hyperlinear. The same holds for
$\EL_n\bigl(\LC(X,\F_2)\rtimes_T\Z\bigr)$ for every $n\ge3$.
\end{theorem}

Property~\textup{(T)} follows from the theorem of Ershov and
Jaikin-Zapirain~\cite{EJZ}. The finite models use periodic approximation,
as in the proof by Grigorchuk and Medynets that topological full groups
of minimal Cantor systems are LEF~\cite[Theorem~2.6]{GM}. The new step is
simplicity. A nontrivial normal subgroup contains a nontrivial commutator
lying in a finite simple group $\GL_d(\F_2)$ acting on a clopen tower, so
it contains this group and with it an elementary matrix. For derived
topological full groups, Matui showed in the same way that a nontrivial
normal subgroup meets a simple union of alternating groups on
towers~\cite[Lemma~3.4 and Theorem~4.9]{Matui}. Thom constructed a
finitely generated Kazhdan LEF group that is not residually
finite~\cite{Thom}, but his example is not simple.

The proof needs only a free minimal subshift whose patterns on each
finite window are realized by a finite subshift, and every residually
finite group acts on such a subshift. So every finitely generated
residually finite group is a subgroup of an infinite simple Kazhdan LEF
group (Corollary~\ref{cor:rf}). The word problem of $G_X$ has the Turing
degree of the language of $X$, and every degree occurs
(Corollary~\ref{cor:wp}).

\section{Proof of Theorem~\ref{thm:main}}

We write the proof for $n=3$. For $n\ge3$ replace $3$ by $n$ throughout,
so that $d=n(2w+1)$ below.

\subsection*{The ring and property (T)}
Write $e_U$ for the indicator of a clopen set $U\subseteq X$. Our
conventions are $(Tx)_t=x_{t+1}$ and
\[
 R=\LC(X,\F_2)\rtimes_T\Z=\bigoplus_{j\in\Z}\LC(X,\F_2)\,u^j,
 \qquad ufu^{-1}=f\circ T^{-1}.
\]
So every element of $R$ is uniquely a finite sum $\sum_jf_ju^j$, and
$ue_Uu^{-1}=e_{TU}$.

The ring $R$ is generated by $u^{\pm1}$ and the letter indicators
$e_a=e_{\{x:x_0=a\}}$, since products of the translates $u^je_au^{-j}$
are the cylinder indicators, which span $\LC(X,\F_2)$. Let $G=\EL_3(R)$
be the subgroup of $\GL_3(R)$ generated by the elementary matrices
$e_{ij}(r)=I_3+rE_{ij}$, $i\ne j$, $r\in R$, where $E_{ij}$ is the usual
matrix unit. With $[g,h]=ghg^{-1}h^{-1}$,
\begin{equation}\label{eq:elementary}
 \begin{aligned}
 e_{ij}(r+s)&=e_{ij}(r)e_{ij}(s),\qquad i\ne j,\\
 [e_{ik}(r),e_{kj}(s)]&=e_{ij}(rs),\qquad i,j,k\text{ distinct}.
 \end{aligned}
\end{equation}
So the matrices $e_{ij}(s)$ with $s\in\{u,u^{-1}\}\cup\{e_a:a\in A\}$
generate $G$. By Ershov and Jaikin-Zapirain~\cite[Theorem~1.1]{EJZ},
$\EL_n(S)$ has property~\textup{(T)} for every finitely generated
associative unital ring $S$ and every $n\ge3$. So $G$ has
property~\textup{(T)}. It is infinite because $e_{12}(\LC(X,\F_2))$ is
infinite.

\subsection*{Finite models}
A group is LEF~\cite{VershikGordon} if every finite subset embeds
injectively into a finite group, preserving all products that stay in
that subset.

Fix $x\in X$ and $\ell\ge0$. By minimality every word of $X$ occurs in
$x$, and $x_{[0,2\ell)}$ recurs at arbitrarily large positions. Choose
such a position $m_\ell$ with all words of length $2\ell+1$ of $X$
occurring in $x_{[0,m_\ell+2\ell)}$. This word begins and ends with
$x_{[0,2\ell)}$, so it has period $m_\ell$, and the $m_\ell$-periodic
sequence $y_\ell$ that agrees with $x$ on $[0,m_\ell)$ has the same words
of length $2\ell+1$ as $X$. So $m_\ell$ is at least the number of these
words, which tends to infinity because $X$ is infinite.

On $\F_2^{\Z/m_\ell\Z}$ let $P\delta_t=\delta_{t+1}$. For $f$ depending
only on coordinates in $[-\ell,\ell]$ let
$D_\ell(f)\delta_t=f(T^ty_\ell)\delta_t$, where $f(T^ty_\ell)$ is the value
of $f$ at any point of $X$ that agrees with $T^ty_\ell$ on $[-\ell,\ell]$,
and put $\varphi_\ell(\sum_jf_ju^j)=\sum_jD_\ell(f_j)P^j$. Since
$PD_\ell(f)P^{-1}=D_\ell(f\circ T^{-1})$, for all $r,s\in R$ the identities
$\varphi_\ell(r+s)=\varphi_\ell(r)+\varphi_\ell(s)$ and
$\varphi_\ell(rs)=\varphi_\ell(r)\varphi_\ell(s)$ hold for large $\ell$. If
$r=\sum_jf_ju^j\ne0$, then for large $\ell$ some $D_\ell(f_j)\ne0$, and the
terms with $|j|<m_\ell/2$ have disjoint supports, so $\varphi_\ell(r)\ne0$.
For a finite $F\subseteq G$ and large $\ell$, applying $\varphi_\ell$ to the
entries of the elements of $F$ and of their inverses gives an injective
map $F\to\GL_{3m_\ell}(\F_2)$ that preserves the products staying in $F$.
So $G$ is LEF. LEF groups are sofic, and sofic groups are
hyperlinear~\cite[Example~4.5 and Theorem~3.3]{Pestov}.

So $L(G)$ embeds in $R^\omega$~\cite[Proposition~7.1]{Ozawa}. Since $G$
is infinite and simple, its nontrivial conjugacy classes are infinite,
so $L(G)\mathbin{\bar\otimes}R$ is a McDuff factor. It embeds in
$R^\omega$, and its unitary group contains $G$, as Brown asked. Kirchberg
proved that a Kazhdan group with the factorization property is
residually finite~\cite[Theorem~1.1]{Kirchberg}. So $G$ does not have the
factorization property, and $C^*(G)$ does not have the local lifting
property~\cite[p.~527]{Ozawa}.

\subsection*{Simplicity}
Let $1\ne N\trianglelefteq G$ and $1\ne g\in N$, and let $w\ge0$ bound
the absolute values of the exponents in the entries of $g$ and $g^{-1}$.
Call a clopen set $V$ small if $V\cap T^jV=\varnothing$ for
$0<|j|\le2w$ and every $f\circ T^i$, with $|i|\le w$ and $f$ a
coefficient of an entry of $g$ or $g^{-1}$, is constant on $V$.
Minimality and infiniteness imply that $T$ has no periodic points, so
every point has a small clopen neighborhood. By compactness every clopen
set is a finite disjoint union of small ones.

Some $h=e_{ij}(e_V)$ with $V$ small does not commute with $g$.
Otherwise $g$ commutes with $e_{ij}(e_V)$ for every clopen $V$,
by~\eqref{eq:elementary}. Taking $V=X$ gives $g=cI_3$ with
$c=\sum_jc_ju^j\in R$. Then $e_Vc-ce_V=\sum_jc_j(e_V-e_{T^jV})u^j$
vanishes for every clopen $V$. Choosing $V$ to contain $x$ but not
$T^{-j}x$, we get $c_j(x)=0$ for $j\ne0$. So $c\in\LC(X,\F_2)$, and
comparing constant coefficients in $cc^{-1}=1$ gives $c=1$ and $g=1$.

Fix such $h$ and $V$, put $d=3(2w+1)$, and for $|a|,|b|\le w$ put
$\epsilon_{ab}=e_{T^aV}u^{a-b}$. Since the sets $T^aV$ are disjoint,
$\epsilon_{ab}\epsilon_{a'b'}=\delta_{ba'}\epsilon_{ab'}$. Index the
coordinates of $\F_2^d$ by pairs $(p,a)$ with $1\le p\le3$ and
$|a|\le w$. Then $\psi(E_{(p,a),(q,b)})=\epsilon_{ab}E_{pq}$ defines an
injective multiplicative linear map $\psi:M_d(\F_2)\to M_3(R)$, and
$A\mapsto I_3-\psi(I_d)+\psi(A)$ embeds $\GL_d(\F_2)$ in $\GL_3(R)$. Its
image $H$ lies in $G$. Indeed, transvections generate $\GL_d(\F_2)$. The
transvection between $(p,a)$ and $(q,b)$ maps to $e_{pq}(\epsilon_{ab})$
if $p\ne q$. If $p=q$, it is the commutator of the transvections between
$(p,a)$ and $(p',a)$ and between $(p',a)$ and $(p,b)$, where $p'\ne p$.

Put $k=[g,h]\in N\setminus\{1\}$. If $|a|,|b|\le w$ and $f,f'$ are
coefficients as above, then
\[
  fu^a\,e_V\,f'u^b=f\,e_{T^aV}\,(f'\circ T^{-a})\,u^{a+b}
  \in\{0,\epsilon_{a,-b}\},
\]
since $f\circ T^a$ and $f'$ are constant on $V$. The entries of
$ghg^{-1}-I_3=g\,e_VE_{ij}\,g^{-1}$ are sums of such products, and
$h^{-1}=h=I_3+\epsilon_{00}E_{ij}$. So $k-I_3=(ghg^{-1}-h)h$ and
$k^{-1}-I_3=h(ghg^{-1}-h)$ lie in $\psi(M_d(\F_2))$. As $\psi$ is
injective and multiplicative, $k\in H$. Then $N\cap H$ is a nontrivial
normal subgroup of $H\cong\GL_d(\F_2)=\PSL_d(\F_2)$, which is simple as
$d\ge3$. So $H\subseteq N$, and $e_{pq}(e_V)\in N$ for all $p\ne q$.

Finally, the level $J=\{r\in R:e_{pq}(r)\in N\text{ for all }p\ne q\}$
of $N$~\cite{Stepanov} is a two-sided ideal. It is additive
by~\eqref{eq:elementary}, and for $l\notin\{p,q\}$ we have
$e_{pq}(sr)=[e_{pl}(s),e_{lq}(r)]$ and $e_{pq}(rs)=[e_{pl}(r),e_{lq}(s)]$.
It contains $e_V$, so it contains every $e_{T^aV}=u^ae_Vu^{-a}$. By
minimality and compactness finitely many of these sets cover $X$, so
$1=1-\prod_i(1-e_{T^{a_i}V})\in J$. So $J=R$ and $N=G$.\hfill$\square$

\section{Residually finite subgroups}

\begin{corollary}\label{cor:rf}
Every finitely generated residually finite group is a subgroup of an
infinite, finitely generated, simple group with property~\textup{(T)}
that is LEF.
\end{corollary}

Without property~\textup{(T)}, Kionke and Schesler embedded every
finitely generated residually finite group in a finitely generated
simple LEF group~\cite[Theorem~1.2]{KionkeSchesler}.

\begin{proof}
Let $\Gamma$ be generated by a finite set $S$. Replacing $\Gamma$ by
$\Gamma\times\Z$, we may assume that $\Gamma$ is infinite. Let
$\Gamma_n$ be normal subgroups of finite index with
$\bigcap_n\Gamma_n=1$. Left multiplication on
$\Gamma/\Gamma_n\times\{1,2\}$ gives even permutations, so $\Gamma$
embeds in $P=\prod_n\operatorname{Sym}(\Gamma/\Gamma_n\times\{1,2\})$.
Every even permutation of a finite set is a commutator of two
permutations~\cite[Theorem~1]{Ore}, so each $s\in S$ equals $[a_s,b_s]$
for some $a_s,b_s\in P$. The subgroup $\Delta$ of $P$ generated by $S$
and the $a_s,b_s$ is finitely generated, infinite and residually finite,
and $\Gamma\le[\Delta,\Delta]$.

Choose normal subgroups $\Delta=K_0>K_1>\cdots$ of finite index with
$\bigcap_mK_m=1$ and $[K_m:K_{m+1}]\ge4$. Define $x\in\{0,1,2\}^\Delta$
by $x(1)=1$ and, on each set $K_m\setminus K_{m+1}$, by $x=1$ on one
coset of $K_{m+1}$, $x=2$ on another, and $x=0$ on the others. Then
$x(hk)=x(h)$ whenever $h\notin K_m$ and $k\in K_m$. Let $\Delta$ act by
$(g\cdot y)(h)=y(hg)$, and let $X$ be the closure of $\Delta\cdot x$.
Arrays of this kind go back to Toeplitz sequences; see Cortez and
Petite~\cite{CortezPetite} for residually finite groups.

The subshift $X$ is minimal. Let $F\subseteq\Delta$ be finite, let
$\sigma\in\Delta$, and choose $m$ with $F\sigma\cap K_m\subseteq\{1\}$.
If $1\notin F\sigma$, the pattern $h\mapsto x(h\sigma)$ on $F$ recurs at
$\sigma k$ for all $k\in K_m$. If $1\in F\sigma$, it recurs at $\sigma k$
for all $k$ in the coset of $K_{m+1}$ in $K_m$ on which $x=1$. So every
pattern of $x$ recurs along a coset of a subgroup of finite index.

The action on $X$ is free. Let $\gamma\ne1$, choose $n$ with
$\gamma\notin K_n$, and let $T$ be a transversal of $K_n$. Fix $\sigma$
and put $\delta=\sigma^{-1}\gamma\sigma$, so $\delta\in K_m\setminus K_{m+1}$
for some $m<n$. The cosets of $K_{m+1}$ in $K_m$ on which $x=1$ and $x=2$
contain points $t\sigma$ and $t'\sigma$ with $t,t'\in T$. Right
multiplication by $\delta$ moves each of them to another coset of
$K_{m+1}$ in $K_m$, and not both land in $K_{m+1}$. So
$x(t\sigma)\ne x(t\sigma\delta)$ or $x(t'\sigma)\ne x(t'\sigma\delta)$.
Since $t\sigma\delta=t\gamma\sigma$, the configurations $\sigma\cdot x$ and
$\gamma\cdot(\sigma\cdot x)$ differ on $T$ for every $\sigma$, so $\gamma$
fixes no point of $X$.

The subshift $X$ has finite models. Let $F\subseteq\Delta$ be finite and
choose $n$ with $FF^{-1}\cap K_n=\{1\}$. For $c\in\{0,1,2\}$ let $p_c$
agree with $x$ off $K_n$ and take the value $c$ on $K_n$. Then
$k\cdot p_c=p_c$ for $k\in K_n$, and $p_0,p_1,p_2$ have together the same
patterns on $F$ as $X$. Indeed, for each $\sigma$ at most one $h\in F$ has
$h\sigma\in K_n$, and $k\mapsto x(h\sigma k)$ takes all three values on
$K_n$.

The proof of Theorem~\ref{thm:main} now applies to
$R=\bigoplus_{g\in\Delta}\LC(X,\F_2)u_g$, with $u_gfu_g^{-1}=f\circ g^{-1}$.
Exponents become group elements, and bounds on exponents become bounds
on word length. For a large ball $F$, the finite model acts on
$\bigoplus_c\F_2^{\Delta/K_n}$ by $u_g\delta_{hK_n}=\delta_{ghK_n}$ and by
$f\delta_{hK_n}=f(h\cdot p_c)\delta_{hK_n}$ on the summand of $c$, and
$FF^{-1}\cap K_n=\{1\}$ separates the terms of distinct $u_g$. In the
simplicity proof, small sets satisfy $V\cap gV=\varnothing$ for $g\ne1$
of length at most $2w$, the tower units are $\epsilon_{ab}=e_{aV}u_{ab^{-1}}$
for $a,b$ of length at most $w$, freeness gives the scalar step, and
finitely many translates $e_{aV}=u_ae_Vu_a^{-1}$ cover $X$. So
$G=\EL_3(R)$ is an infinite, finitely generated, simple group with
property~\textup{(T)} that is LEF.

Over $\F_2$, for a unit $c$ of $R$,
\[
 e_{12}(c)\,e_{21}(c^{-1})\,e_{12}(c)\;e_{12}(1)\,e_{21}(1)\,e_{12}(1)
 =\diag(c,c^{-1},1).
\]
So for units $a,b$ the matrix
\[
 \diag(a,a^{-1},1)\,\diag(b,b^{-1},1)\,\diag\bigl((ba)^{-1},ba,1\bigr)
 =\diag([a,b],1,1)
\]
lies in $G$. Since $\Gamma\le[\Delta,\Delta]$, the map
$\gamma\mapsto\diag(u_\gamma,1,1)$ is an injective homomorphism
$\Gamma\to G$.
\end{proof}

\section{Word problems}

\begin{corollary}\label{cor:wp}
The word problem of $G_X$ has the Turing degree of the language $L(X)$
of $X$. Every Turing degree occurs, so there are continuum many pairwise
nonisomorphic groups $G_X$.
\end{corollary}

\begin{proof}
Multiplying out a word in the generators gives a matrix with entries
$\sum_jf_ju^j$, each $f_j$ a table on words, using $uf=(f\circ T^{-1})u$.
The word is trivial if and only if the tables of its difference from
$I_3$ vanish on $L(X)$, so $L(X)$ computes the word problem. Conversely,
by~\eqref{eq:elementary} a word for $e_{12}(\prod_{t<n}u^{-t}e_{v_t}u^t)$
is computable from $v=v_0\cdots v_{n-1}$, and it is trivial if and only
if $v\notin L(X)$. For derived topological full groups, Grigorchuk and
Medynets proved that the word problem is decidable if and only if
$L(X)$ is recursive~\cite[Theorem~1.1(3)]{GMpres}.

For irrational $\alpha\in(0,1)$ let $X_\alpha$ be the infinite minimal
Sturmian subshift, the closure of the codings $c(\theta)$,
$\theta\in[0,1)$, with $c(\theta)_t=1$ if and only if
$\theta+t\alpha\bmod1\ge1-\alpha$~\cite{MorseHedlund}. Its words of
length $n$ are the constant values of $c(\theta)_{[0,n)}$ on the arcs
between the points $-j\alpha\bmod1$, $0\le j\le n$, so $\alpha$ computes
$L(X_\alpha)$. The word $c(\theta)_{[0,n)}$ has
$\lfloor\theta+n\alpha\rfloor$ ones, within $1$ of $n\alpha$, so
$L(X_\alpha)$ computes $\alpha$. Every set $S\subseteq\mathbb N$ has the
degree of the irrational number $[0;1+\chi_S(0),1+\chi_S(1),\dots]$, and
the degree of the word problem is an isomorphism invariant.
\end{proof}

\section*{Questions}
A finitely presented LEF group is residually finite~\cite{VershikGordon},
and an infinite simple group is not, so no group $G_X$ is finitely
presented.
\begin{enumerate}
\item Is some finitely presented infinite simple group with
property~\textup{(T)} sofic, or at least hyperlinear? A positive answer
would also answer Open problem~6.1 of Alekseev and Thom~\cite{AlekseevThom}.
\item Does every finitely generated LEF group embed in an infinite simple
Kazhdan LEF group?
\item Flip conjugate subshifts give isomorphic groups. Does
$G_X\cong G_Y$ imply that $X$ and $Y$ are flip conjugate, or at least
strongly orbit equivalent~\cite{GPS}?
\end{enumerate}

\subsection*{Origin and authorship}
Claude (Anthropic) found the construction and original proofs under
the author's direction on September~12, 2026, and wrote the original
manuscript. Codex (OpenAI) shortened an earlier version, and Claude
revised this one. The author is responsible for the final manuscript.
Proof records are in the
\href{https://github.com/SauersML/group-approximation}{project repository}.

\begin{thebibliography}{99}

\bibitem{AlekseevThom}
V.~Alekseev and A.~Thom,
\emph{Centralizers of sofic approximations of Kazhdan groups},
\href{https://arxiv.org/abs/2608.05362}{arXiv:2608.05362} (2026).

\bibitem{Brown}
N.~P. Brown, \emph{Tracial invariants, classification and
$\mathrm{II}_1$ factor representations of Popa algebras},
\href{https://arxiv.org/abs/math/0111286}{arXiv:math/0111286} (2001).

\bibitem{CortezPetite}
M.~I. Cortez and S.~Petite,
\emph{$G$-odometers and their almost one-to-one extensions},
J. Lond. Math. Soc. (2) \textbf{78} (2008), 1--20.
\doi{10.1112/jlms/jdn002}.

\bibitem{EJZ}
M.~Ershov and A.~Jaikin-Zapirain,
\emph{Property~\textup{(T)} for noncommutative universal lattices},
Invent. Math. \textbf{179} (2010), 303--347.
\doi{10.1007/s00222-009-0218-2}.

\bibitem{GPS}
T.~Giordano, I.~F. Putnam, and C.~F. Skau,
\emph{Topological orbit equivalence and $C^*$-crossed products},
J. Reine Angew. Math. \textbf{469} (1995), 51--112.
\doi{10.1515/crll.1995.469.51}.

\bibitem{GM}
R.~Grigorchuk and K.~Medynets,
\emph{On algebraic properties of topological full groups},
Sb. Math. \textbf{205} (2014), 843--861.
\doi{10.1070/SM2014v205n06ABEH004400}.

\bibitem{GMpres}
R.~Grigorchuk and K.~Medynets,
\emph{Presentations of topological full groups by generators and
relations}, J. Algebra \textbf{500} (2018), 46--68.
\doi{10.1016/j.jalgebra.2016.10.027}.

\bibitem{KionkeSchesler}
S.~Kionke and E.~Schesler,
\emph{From telescopes to frames and simple groups},
J. Comb. Algebra (2024), \doi{10.4171/jca/103};
\href{https://arxiv.org/abs/2304.09307}{arXiv:2304.09307}.

\bibitem{Kirchberg}
E.~Kirchberg, \emph{Discrete groups with Kazhdan's property~\textup{T}
and factorization property are residually finite},
Math. Ann. \textbf{299} (1994), 551--563.
\doi{10.1007/BF01459798}.

\bibitem{Matui}
H.~Matui, \emph{Some remarks on topological full groups of Cantor
minimal systems}, Internat. J. Math. \textbf{17} (2006), 231--251.
\doi{10.1142/S0129167X06003448}.

\bibitem{MorseHedlund}
M.~Morse and G.~A. Hedlund,
\emph{Symbolic dynamics II. Sturmian trajectories},
Amer. J. Math. \textbf{62} (1940), 1--42.
\doi{10.2307/2371431}.

\bibitem{Ore}
O.~Ore, \emph{Some remarks on commutators},
Proc. Amer. Math. Soc. \textbf{2} (1951), 307--314.
\doi{10.1090/S0002-9939-1951-0040298-4}.

\bibitem{Ozawa}
N.~Ozawa, \emph{About the QWEP conjecture},
Internat. J. Math. \textbf{15} (2004), 501--530.
\doi{10.1142/S0129167X04002417}.

\bibitem{Pestov}
V.~G. Pestov, \emph{Hyperlinear and sofic groups: a brief guide},
Bull. Symbolic Logic \textbf{14} (2008), 449--480.
\doi{10.2178/bsl/1231081461}.

\bibitem{Stepanov}
A.~V. Stepanov, \emph{On the normal structure of the general linear
group over a ring}, Zap. Nauchn. Sem. POMI \textbf{236} (1997), 166--182;
English transl., J. Math. Sci. \textbf{95} (1999), 2146--2155.
\doi{10.1007/BF02169976}.

\bibitem{Thom}
A.~Thom, \emph{Examples of hyperlinear groups without factorization
property}, Groups Geom. Dyn. \textbf{4} (2010), 195--208.
\doi{10.4171/GGD/80}.

\bibitem{VershikGordon}
A.~M. Vershik and E.~I. Gordon,
\emph{Groups that are locally embeddable in the class of finite groups},
Algebra i Analiz \textbf{9} (1997), no.~1, 71--97; English transl.,
St. Petersburg Math. J. \textbf{9} (1998), no.~1, 49--67.

\end{thebibliography}
\end{document}
